\section[Conclusão]{Conclusão}

Ao avaliar o resultado do Seniors -- Empregabilidade, considero que entregamos um
MVP muito satisfatório. Todas as funcionalidades mantidas no escopo final foram
implementadas, e a retirada do envio de arquivos de currículo foi acordada com as
\textit{stakeholders}. Como uma versão 0.1, a solução atingiu bem seu objetivo,
percepção reforçada pela satisfação demonstrada na apresentação final. A
infraestrutura também permaneceu estável durante praticamente todo o semestre, com
grande mérito do colega AGES III Leonardo Dondoni. A ausência de um domínio próprio
e o uso dos e-mails padrão do Cognito foram limitações aceitáveis diante do tempo e
do orçamento disponíveis.

Apesar da qualidade do produto, não considero este o melhor projeto das minhas
quatro experiências na AGES. O resultado foi mais completo e mais bem recebido do
que os projetos das minhas passagens por AGES I e AGES III, mas o desenvolvimento
foi difícil. A presença de cinco alunos AGES IV e cinco AGES III tornou mais
trabalhosas as decisões que deveriam orientar o time. Ainda assim, guardo uma
lembrança positiva do grupo: as aulas tiveram brincadeiras, conversas e um ambiente
descontraído, além da orientação da Prof.ª Dra. Cristina Moreira Nunes, por quem
tenho grande consideração.

Minha principal contribuição foi a comunicação. Procurei iniciar conversas, reunir
o time para as reuniões diárias, cobrar a atualização das tarefas e manter as
pessoas próximas. Também valorizava os momentos informais e até o hábito de levar
chocolates para as aulas, pois acredito que um ambiente leve facilita pedir ajuda e
trabalhar em conjunto. Minha maior falha, por outro lado, foi não organizar o backlog
adequadamente desde a Sprint 0. Histórias, épicos e tarefas não tiveram a
granularidade necessária, e o ClickUp se tornou difícil de acompanhar. As
visualizações alternativas e os filtros por squad reduziram o problema, mas não o
resolveram. Eu poderia ter defendido com mais firmeza uma organização granular no
Notion ou exigido que a estrutura do backlog fosse validada antes do
desenvolvimento. Percebi que tentar agradar a todos pode adiar decisões necessárias
e prejudicar a própria equipe.

A divisão por squads permaneceu relativamente estável, mas esperávamos que os
alunos AGES III atuassem mais como líderes técnicos, orientando os iniciantes e
revisando os PRs. Na minha avaliação, muitos assumiram tarefas individuais, e nós,
alunos AGES IV, concentramos parte da cobrança, da revisão e da integração. Essa
crítica não se aplica igualmente a todas as contribuições, como demonstra o trabalho
do Leonardo na infraestrutura, mas foi uma dificuldade recorrente. Também senti uma
colaboração mais individual do que em projetos anteriores. O uso mais presente de
ferramentas de inteligência artificial pode ter influenciado essa dinâmica, embora
não considere que tenha sido sua única causa.

Os conhecimentos das disciplinas do curso foram necessários em todas as etapas:
requisitos orientaram as histórias de usuário; a modelagem de dados sustentou
currículos, vagas e candidaturas; a arquitetura permitiu disponibilizar a aplicação;
e os testes foram essenciais antes da entrega. A AGES acrescentou o desafio de
coordenar pessoas com diferentes níveis de experiência para que essas atividades
acontecessem no momento adequado. Nas primeiras sprints, atuamos de forma reativa e
dependemos de mobilizações próximas ao \textit{code freeze}. Na última, encerramos o
desenvolvimento com antecedência, testamos o sistema em produção e preparamos a
demonstração com calma. Isso mostrou que a pressão anterior estava relacionada à
integração tardia e à pouca visibilidade do trabalho.

Encerro a AGES IV satisfeito com o produto e com a evolução do time, mas consciente
de que um bom resultado não torna o processo automaticamente bom. Aprendi que gerir
exige estruturar o trabalho, tornar bloqueios visíveis, acompanhar as integrações e
tomar decisões mesmo quando não será possível agradar a todos. Ao mesmo tempo,
continuo acreditando que a comunicação e um ambiente acolhedor são fundamentais. O
equilíbrio entre firmeza e proximidade foi a principal lição que esta experiência
deixou para minha atuação profissional.
